\begin{abstract}
\noindent
Quanto do mercado de trabalho brasileiro está efetivamente sujeito à inteligência artificial --- e quem são os expostos? Este artigo desenvolve um modelo de designação de tarefas com adoção corporativa e individual de IA sob informalidade, e testa suas previsões nos microdados da PNAD Contínua, combinados ao índice de \citet{eloundou2024gpts} por um cruzamento ocupacional determinístico. A exposição concentra-se no núcleo formal e escolarizado: 35\% da força de trabalho (35{,}7 milhões de pessoas) atua em ocupações de exposição quase nula ($\theta \le 2$). Em modelos no nível do indivíduo com controles mincerianos completos (idade, gênero e raça), cada ano adicional de estudo associa-se a 0{,}23 ponto a mais de exposição à IA (erro-padrão 0{,}03 agrupado em 122 ocupações; $N = 227.629$). A associação entre exposição e formalidade é robusta, e o gradiente é mais íngreme nas 27 Unidades da Federação com setor formal mais profundo. No Brasil, a informalidade funciona como o principal amortecedor da exposição à IA.
\end{abstract}
